\documentclass[12pt]{article}
\usepackage[utf8]{inputenc}
\usepackage[T1]{fontenc}
\usepackage{hyperref}
\title{Kajali Language Resources for K-12 Curriculum Development}
\author{Compiled by Claude Sonnet 5.5 (claude-sonnet-5-5) for Andr\'as Kornai}
\date{}
\begin{document}
\maketitle

\section{Introduction}
Kajali (ISO 639-3 xkj; Glottocode kaja1247) is a Tati dialect or language of northwestern Iran, spoken in the village of Kajal in Khalkhal County, Ardabil Province. Glottolog places it in the Khalkhalic branch of Central Tatic. The 2016 census counted 227 residents of Kajal. Sources disagree on its vitality, and one calls it moribund. It is mainly spoken rather than written. The material found is purely descriptive, and no teaching materials were found.

\section{Primary Educational Resources}
None found. Kajali is not known to be taught in schools.

\section{Linguistic and Reference Materials}
\begin{itemize}
\item Wikipedia, Kajali language: \url{https://en.wikipedia.org/wiki/Kajali_language}
\item Wikipedia, Kajal, Iran: \url{https://en.wikipedia.org/wiki/Kajal,_Iran}
\item Wikipedia, Tati language (Iran): \url{https://en.wikipedia.org/wiki/Tati_language_(Iran)}
\item Glottolog entry: \url{https://glottolog.org/resource/languoid/id/kaja1247.bigmap.html}
\item Bulletin of the School of Oriental and African Studies, ``The Tati dialect of Kajal'' (the title is as shown in the search result; the content was not read): \url{https://resolve.cambridge.org/core/journals/bulletin-of-the-school-of-oriental-and-african-studies/article/tati-dialect-of-kajal/7EFEFA45AB2E3362C0C5456D8A674661}
\item Verbix endangered-languages page: \url{https://docs.verbix.com/EndangeredLanguages/Kajali}
\end{itemize}

\section{Religious and Community Resources}
None found.

\section{Digital and Audio Resources}
None found.

\section{Curriculum Development Context}
With a very small speaker base in one village and no known written tradition, the realistic first step is documentation: recordings, a wordlist, and an orthography, based on Persian script. The existing academic description of the dialect could be a starting point. Further work would need local community involvement.
\end{document}